\begin{ex}
Sia $A:=\{0,1\}$ l'alfabeto binario. Stabilire se, per ogni dato $n$ intero positivo, l'insieme $X_n\subset A^*$
\[
X_n:=\{a^ib^jc^k \mid i+j\geq 3,\ k\equiv 0 \bmod n \text{ e } a,b,c\in A\}
\]
è DFA-decidibile. In caso affermativo fornire un automa che lo decide.
\end{ex}

\begin{sol}
L'insieme $X_n$ è DFA-decidibile per ogni $n$. Per decidere il primo pezzo, cioè $a^ib^j$ con $i+j\geq 3$, basta l'automa $\mathcal{A}$ qui sotto (le transizioni non disegnate vanno in uno stato pozzo).
\[
\begin{tikzpicture}[
    scale=0.8, transform shape,
    >=stealth', shorten >=1pt, auto,
    every state/.style={circle, draw, fill=white, minimum size=0.8cm, inner sep=1pt}
]
% Stati
\node[state] (qe)    at (0,0)     {$q_\varepsilon$};
\node[state] (q0)    at (2,2.75)  {$q_0$};
\node[state] (q1)    at (2,-2.75) {$q_1$};
\node[state] (q00)   at (4,4)     {$q_{00}$};
\node[state] (q01)   at (4,1.5)   {$q_{01}$};
\node[state] (q10)   at (4,-1.5)  {$q_{10}$};
\node[state] (q11)   at (4,-4)    {$q_{11}$};
\node[state] (q000)  at (6.5,5)   {$q_{000}$};
\node[state] (q001)  at (6.5,3)   {$q_{001}$};
\node[state] (q011)  at (6.5,1)   {$q_{011}$};
\node[state] (q100)  at (6.5,-1)  {$q_{100}$};
\node[state] (q110)  at (6.5,-3)  {$q_{110}$};
\node[state] (q111)  at (6.5,-5)  {$q_{111}$};
\node[state] (q0001) at (9.5,5)   {$q_{0001}$};
\node[state] (q1110) at (9.5,-5)  {$q_{1110}$};
\node[state, accepting] (qF) at (12,0) {$q_F$};

% Freccia d'ingresso
\draw[->] ([shift={(-0.7,0.7)}]qe.center) -- (qe);

% Transizioni
\draw[->] (qe)  -- node {$0$} (q0);
\draw[->] (qe)  -- node[swap] {$1$} (q1);
\draw[->] (q0)  -- node {$0$} (q00);
\draw[->] (q0)  -- node[swap] {$1$} (q01);
\draw[->] (q1)  -- node {$0$} (q10);
\draw[->] (q1)  -- node[swap] {$1$} (q11);
\draw[->] (q00) -- node {$0$} (q000);
\draw[->] (q00) -- node[swap] {$1$} (q001);
\draw[->] (q01) -- node[swap] {$1$} (q011);
\draw[->] (q10) -- node {$0$} (q100);
\draw[->] (q11) -- node {$0$} (q110);
\draw[->] (q11) -- node[swap] {$1$} (q111);

% Passaggi 000 -> 0001 e 111 -> 1110
\draw[->] (q000) -- node {$1$} (q0001);
\draw[->] (q111) -- node[swap] {$0$} (q1110);

% Cappi
\draw[->] (q000)  to[loop above, looseness=8] node {$0$} (q000);
\draw[->] (q001)  to[loop above, looseness=8] node {$1$} (q001);
\draw[->] (q011)  to[loop above, looseness=8] node {$1$} (q011);
\draw[->] (q100)  to[loop below, looseness=8] node {$0$} (q100);
\draw[->] (q110)  to[loop below, looseness=8] node {$0$} (q110);
\draw[->] (q111)  to[loop below, looseness=8] node {$1$} (q111);
\draw[->] (q0001) to[loop above, looseness=8] node {$1$} (q0001);
\draw[->] (q1110) to[loop below, looseness=8] node {$0$} (q1110);

% Epsilon-mosse verso lo stato finale
\draw[->] (q000)  to[bend left=8]  node[pos=0.7] {$\varepsilon$} (qF);
\draw[->] (q001)  -- node[pos=0.7] {$\varepsilon$} (qF);
\draw[->] (q011)  -- node[pos=0.7] {$\varepsilon$} (qF);
\draw[->] (q100)  -- node[pos=0.7] {$\varepsilon$} (qF);
\draw[->] (q110)  -- node[pos=0.7] {$\varepsilon$} (qF);
\draw[->] (q111)  to[bend right=8] node[pos=0.7] {$\varepsilon$} (qF);
\draw[->] (q0001) -- node {$\varepsilon$} (qF);
\draw[->] (q1110) -- node[swap] {$\varepsilon$} (qF);
\end{tikzpicture}
\]
Lo concateniamo con l'automa $\mathcal{B}_n$, che accetta $\{c^k \mid k\equiv 0 \bmod n\}$. Per $n=1$ si ha $r_{0^1}=r_{0^n}$ e $r_{1^1}=r_{1^n}$.
\[
\begin{tikzpicture}[
    scale=0.8, transform shape,
    >=stealth', shorten >=1pt, auto,
    every state/.style={circle, draw, fill=white, minimum size=0.8cm, inner sep=1pt}
]
% Stati
\node[state, accepting] (re)  at (0,0)     {$r_\varepsilon$};
\node[state]            (r0)  at (3,1.5)   {$r_0$};
\node[state]            (r02) at (6,1.5)   {$r_{0^2}$};
\node                   (r0d) at (8.5,1.5) {$\cdots$};
\node[state, accepting] (r0n) at (11,1.5)  {$r_{0^n}$};
\node[state]            (r1)  at (3,-1.5)  {$r_1$};
\node[state]            (r12) at (6,-1.5)  {$r_{1^2}$};
\node                   (r1d) at (8.5,-1.5){$\cdots$};
\node[state, accepting] (r1n) at (11,-1.5) {$r_{1^n}$};

% Freccia d'ingresso
\draw[->] ([shift={(-0.7,0.7)}]re.center) -- (re);

% Transizioni
\draw[->] (re)  -- node {$0$} (r0);
\draw[->] (re)  -- node[swap] {$1$} (r1);
\draw[->] (r0)  -- node {$0$} (r02);
\draw[->] (r02) -- node {$0$} (r0d);
\draw[->] (r0d) -- node {$0$} (r0n);
\draw[->] (r1)  -- node[swap] {$1$} (r12);
\draw[->] (r12) -- node[swap] {$1$} (r1d);
\draw[->] (r1d) -- node[swap] {$1$} (r1n);

% Ritorni
\draw[->] (r0n) to[bend right=35] node[above] {$0$} (r0);
\draw[->] (r1n) to[bend left=35]  node[below] {$1$} (r1);
\end{tikzpicture}
\]
L'automa che decide $X_n$ è il seguente: gli stati finali sono quelli di $\mathcal{B}_n$, ossia $r_\varepsilon$, $r_{0^n}$ e $r_{1^n}$.
\[
\begin{tikzpicture}[
    >=stealth', shorten >=1pt, auto,
    every state/.style={circle, draw, fill=white, minimum size=0.8cm, inner sep=1pt}
]
% Rettangolo di A
\draw[-, line width=2.5pt, rounded corners=3pt] (0,-1.4) rectangle (4,1.4);
\node at (2,0) {$\mathcal{A}$};

% Rettangolo di B_n
\draw[-, line width=2.5pt, rounded corners=3pt] (7,-1.4) rectangle (11,1.4);
\node at (9,0) {$\mathcal{B}_n$};

% Stati di A
\node[state] (qe) at (0,0) {$q_\varepsilon$};
\node[state] (qF) at (4,0) {$q_F$};

% Stati di B_n
\node[state, accepting] (re)  at (7,0)     {$r_\varepsilon$};
\node[state, accepting] (r0n) at (11,0.8)  {$r_{0^n}$};
\node[state, accepting] (r1n) at (11,-0.8) {$r_{1^n}$};

% Freccia d'ingresso e epsilon-mossa di collegamento
\draw[->] ([shift={(-0.7,0.7)}]qe.center) -- (qe);
\draw[->] (qF) -- node {$\varepsilon$} (re);
\end{tikzpicture}
\]
Abbiamo così mostrato che
l'iniseme $X_n$ è $\epsilon$-NDFA
decidibile e quindi è anche
$DFA$ decidibile.
\end{sol}

\newpage

\begin{ex}
Con $X_n$ come nell'esercizio precedente,
\begin{enumerate}
    \item dare l'espressione regolare $\mathsf{e}_3$ tale che $[\![\mathsf{e}_3]\!]=X_3$,
    \item generalizzare ad una famiglia di espressioni regolari $\mathsf{e}_n$ tale che $[\![\mathsf{e}_n]\!]=X_n$,
\end{enumerate}
(la doppia parentesi quadra indica l'interpretazione insiemistica delle espressioni regolari).
\end{ex}

\begin{sol}
Poiché $X_n$ è la concatenazione di $\{a^ib^j \mid i+j\geq 3\}$ con $\{c^k \mid k\equiv 0 \bmod n\}$, cerchiamo un'espressione regolare per ciascuno dei due pezzi e le concateniamo.

Leggendo l'automa $\mathcal{A}$ dell'esercizio precedente si ottiene
\[
\mathsf{e}_{\mathcal{A}} := 0\Bigl(0\bigl(00^*1^* + 11^*\bigr) + 111^*\Bigr) + 1\Bigl(000^* + 1\bigl(00^* + 11^*0^*\bigr)\Bigr)
\]
che è l'espressione per $\mathcal{A}$.

Per trovare l'espressione di $\mathcal{B}_n$ facciamo le riduzioni.

Passo 0. Partiamo da $\mathcal{B}_n$.
\[
\begin{tikzpicture}[
    scale=0.8, transform shape,
    >=stealth', shorten >=1pt, auto,
    every state/.style={circle, draw, fill=white, minimum size=0.8cm, inner sep=1pt}
]
% Stati
\node[state, accepting] (re)  at (2.5,0)    {$r_\varepsilon$};
\node[state]            (r0)  at (5.5,1.5)  {$r_0$};
\node[state]            (r02) at (8,1.5)    {$r_{0^2}$};
\node                   (r0d) at (10,1.5)   {$\cdots$};
\node[state, accepting] (r0n) at (12,1.5)   {$r_{0^n}$};
\node[state]            (r1)  at (5.5,-1.5) {$r_1$};
\node[state]            (r12) at (8,-1.5)   {$r_{1^2}$};
\node                   (r1d) at (10,-1.5)  {$\cdots$};
\node[state, accepting] (r1n) at (12,-1.5)  {$r_{1^n}$};

% Freccia d'ingresso
\draw[->] ([shift={(-0.7,0.7)}]re.center) -- (re);

% Transizioni
\draw[->] (re)  -- node {$0$} (r0);
\draw[->] (re)  -- node[swap] {$1$} (r1);
\draw[->] (r0)  -- node {$0$} (r02);
\draw[->] (r02) -- node {$0$} (r0d);
\draw[->] (r0d) -- node {$0$} (r0n);
\draw[->] (r1)  -- node[swap] {$1$} (r12);
\draw[->] (r12) -- node[swap] {$1$} (r1d);
\draw[->] (r1d) -- node[swap] {$1$} (r1n);

% Ritorni
\draw[->] (r0n) to[bend right=35] node[above] {$0$} (r0);
\draw[->] (r1n) to[bend left=35]  node[below] {$1$} (r1);
\end{tikzpicture}
\]

Passo 1. Aggiungiamo un nuovo stato iniziale $r_\varepsilon^*$ e un nuovo stato finale $r_F$. Gli stati che erano finali, cioè $r_\varepsilon$, $r_{0^n}$ e $r_{1^n}$, vanno in $r_F$ con una $\mathbb 1$-mossa.
\[
\begin{tikzpicture}[
    scale=0.8, transform shape,
    >=stealth', shorten >=1pt, auto,
    every state/.style={circle, draw, fill=white, minimum size=0.8cm, inner sep=1pt}
]
% Stati
\node[state]            (rs)  at (0,0)      {$r_\varepsilon^*$};
\node[state]            (re)  at (2.5,0)    {$r_\varepsilon$};
\node[state]            (r0)  at (5.5,1.5)  {$r_0$};
\node[state]            (r02) at (8,1.5)    {$r_{0^2}$};
\node                   (r0d) at (10,1.5)   {$\cdots$};
\node[state]            (r0n) at (12,1.5)   {$r_{0^n}$};
\node[state]            (r1)  at (5.5,-1.5) {$r_1$};
\node[state]            (r12) at (8,-1.5)   {$r_{1^2}$};
\node                   (r1d) at (10,-1.5)  {$\cdots$};
\node[state]            (r1n) at (12,-1.5)  {$r_{1^n}$};
\node[state, accepting] (rF)  at (2.5,-4.2) {$r_F$};

% Freccia d'ingresso e nuove 1-mosse
\draw[->] ([shift={(-0.7,0.7)}]rs.center) -- (rs);
\draw[->, very thick] (rs) -- node {$\mathbb 1$} (re);
\draw[->, very thick] (re) -- node[left] {$\mathbb 1$} (rF);
\draw[->, very thick, rounded corners=8pt]
    (r1n.south) -- (r1n.south |- 0,-4.2) -- node[above] {$\mathbb 1$} (rF.east);
\draw[->, very thick, rounded corners=8pt]
    (r0n.east) -- (13.3,1.5) -- (13.3,-5.5) -- node[below] {$\mathbb 1$} (rF.south |- 0,-5.5) -- (rF.south);

% Transizioni dell'automa originale
\draw[->] (re)  -- node {$0$} (r0);
\draw[->] (re)  -- node[swap] {$1$} (r1);
\draw[->] (r0)  -- node {$0$} (r02);
\draw[->] (r02) -- node {$0$} (r0d);
\draw[->] (r0d) -- node {$0$} (r0n);
\draw[->] (r1)  -- node[swap] {$1$} (r12);
\draw[->] (r12) -- node[swap] {$1$} (r1d);
\draw[->] (r1d) -- node[swap] {$1$} (r1n);
\draw[->] (r0n) to[bend right=35] node[above] {$0$} (r0);
\draw[->] (r1n) to[bend left=35]  node[below] {$1$} (r1);
\end{tikzpicture}
\]

Passo 2. Comprimiamo le due catene: da $r_0$ a $r_{0^n}$ si leggono $n-1$ zeri, da $r_1$ a $r_{1^n}$ si leggono $n-1$ uni.
\[
\begin{tikzpicture}[
    scale=0.8, transform shape,
    >=stealth', shorten >=1pt, auto,
    every state/.style={circle, draw, fill=white, minimum size=0.8cm, inner sep=1pt}
]
% Stati
\node[state]            (rs)  at (0,0)      {$r_\varepsilon^*$};
\node[state]            (re)  at (2.5,0)    {$r_\varepsilon$};
\node[state]            (r0)  at (5.5,1.5)  {$r_0$};
\node[state]            (r0n) at (12,1.5)   {$r_{0^n}$};
\node[state]            (r1)  at (5.5,-1.5) {$r_1$};
\node[state]            (r1n) at (12,-1.5)  {$r_{1^n}$};
\node[state, accepting] (rF)  at (2.5,-4.2) {$r_F$};

% Freccia d'ingresso
\draw[->] ([shift={(-0.7,0.7)}]rs.center) -- (rs);
\draw[->] (rs) -- node {$\mathbb 1$} (re);
\draw[->] (re) -- node[left] {$\mathbb 1$} (rF);
\draw[->, rounded corners=8pt]
    (r1n.south) -- (r1n.south |- 0,-4.2) -- node[above] {$\mathbb 1$} (rF.east);
\draw[->, rounded corners=8pt]
    (r0n.east) -- (13.3,1.5) -- (13.3,-5.5) -- node[below] {$\mathbb 1$} (rF.south |- 0,-5.5) -- (rF.south);

% Transizioni
\draw[->] (re) -- node {$0$} (r0);
\draw[->] (re) -- node[swap] {$1$} (r1);
\draw[->, very thick] (r0) -- node {$0^{n-1}$} (r0n);
\draw[->, very thick] (r1) -- node[swap] {$1^{n-1}$} (r1n);
\draw[->] (r0n) to[bend right=35] node[above] {$0$} (r0);
\draw[->] (r1n) to[bend left=35]  node[below] {$1$} (r1);
\end{tikzpicture}
\]

Passo 3. Eliminiamo $r_0$, $r_{0^n}$, $r_1$, $r_{1^n}$. Da $r_0$ si legge $0^{n-1}$ e poi si esce, oppure si torna in $r_0$ leggendo $0$ e si ripete, quindi da $r_\varepsilon$ a $r_F$ passando per questi stati si ha $0(0^{n-1}0)^*0^{n-1}=0(0^n)^*0^{n-1}$. Lo stesso vale per gli uni.
\[
\begin{tikzpicture}[
    scale=0.8, transform shape,
    >=stealth', shorten >=1pt, auto,
    every state/.style={circle, draw, fill=white, minimum size=0.8cm, inner sep=1pt}
]
% Stati
\node[state]            (rs) at (0,0) {$r_\varepsilon^*$};
\node[state]            (re) at (3,0) {$r_\varepsilon$};
\node[state, accepting] (rF) at (9,0) {$r_F$};

% Freccia d'ingresso
\draw[->] ([shift={(-0.7,0.7)}]rs.center) -- (rs);
\draw[->] (rs) -- node {$\mathbb 1$} (re);
\draw[->] (re) -- node {$\mathbb 1$} (rF);

% Nuove transizioni
\draw[->, very thick] (re) to[bend left=40]  node {$0(0^n)^*0^{n-1}$} (rF);
\draw[->, very thick] (re) to[bend right=40] node[swap] {$1(1^n)^*1^{n-1}$} (rF);
\end{tikzpicture}
\]
Quindi $\mathsf{e}_{\mathcal{B}_n} := \varepsilon + 0(0^n)^*0^{n-1} + 1(1^n)^*1^{n-1}$, che equivale a $(0^n)^*+(1^n)^*$ (per $n=1$ la formula resta valida con $0^0=1^0=\varepsilon$).

\begin{enumerate}
    \item Per $n=3$ si ha $\mathsf{e}_{\mathcal{B}_3}=(000)^*+(111)^*$, quindi
    \[
    \mathsf{e}_3 := \mathsf{e}_{\mathcal{A}}\bigl((000)^*+(111)^*\bigr).
    \]
    \item Per ogni $n$ la concatenazione di $\mathsf{e}_{\mathcal{A}}$ con $\mathsf{e}_{\mathcal{B}_n}$ dà
    \[
    \mathsf{e}_n := \mathsf{e}_{\mathcal{A}}\bigl((0^n)^*+(1^n)^*\bigr),
    \]
    e $[\![\mathsf{e}_n]\!]=[\![\mathsf{e}_{\mathcal{A}}]\!]\cdot[\![\mathsf{e}_{\mathcal{B}_n}]\!]=X_n$.
\end{enumerate}
\end{sol}

\newpage

\begin{ex}
{Sia $X \in \operatorname{DTIME}_1^A(c \cdot n^2)$ e sia $X' \subset X$ DFA-decidibile.}
{Mostrare che l'insieme}
\[
X'' = \{ w \mid ww \in X \text{ e } w \notin X' \} \in \operatorname{DTIME}_1^A(c' \cdot n^2).
\]
{(si precisa che l'esistenza di macchine di Turing eventualmente utilizzate, va argomentata nel dettaglio - ad esempio fornendone una rappresentazione come grafo etichettato).}
\end{ex}
\begin{sol}
    Sia $\alpha$ una macchina du Turing mononastro
    che decide $X$ in tempo $cn^2$ ed $A$
    un DFA che decide $X'$.
    Possiamo associare ad $A$
    una macchina di Turing mononastro
    $\beta$ che decide $X'$ in tempo lineare,
    quindi esiste $b\in \mathbb N$ tale che
    $\beta $ decide $X'$
    in tempo $bn^2$.
    Indichiamo con $\gamma$
    la macchina che data una parola $w$,
    scrive sul nastro la parola $ww$,
    ovvero la macchina nella soluzione dell'esercizio 
    $\ref{ex: macchina copia}$.
    Indicando con $\delta$
    la machcina di Turing che
    effettua passi indietro
    fino a quando non incontra due blank,
    possiamo definire la macchina $\mu$
    come la concatenazione seguente
    \[
    \mu:= \quad \beta
    \to \delta \to \gamma \to \delta \to \alpha 
    \]
    Abbiamo chiaramente che la macchina $\mu$
    decide l'insieme $X''$, dobbiamo ora capire in quanto tempo.
    Consideriamo quindi
    una parola $w$ 
    in input di lungezza 
    $n$:
    \begin{enumerate}
        \item 
    la macchina $\beta$ compie
    al più $bn^2$ ;

    \item la macchina $\beta$,
    non ha modificato 
    la parola e si arresta nel primo 
    blank dopo la parola,
    quindi la macchina $\delta$ compie $n+2$ passi,
    di cui $n+1$ per incontrare il primo blank prima della parola
    e poi un passo per posizionarsi all'inizio della parola;

    \item a questo punto $\gamma$ compie $T(n)$ passi;
    
    \item di nuovo $\delta$ compie $(2n) +2$ passi poiché
    la parola è $ww$ ed è lunga $2n$;

    \item infine $\alpha$ compie $(cn)^2$ passi.
    \end{enumerate}
    Quindi, abbiamo che $\mu$
    decide $X''$ in tempo
    \[
    bn^2 + (n+2) + T(n) + (2n+2) + c^2n^2
    \]
    Mostriamo ora che la macchina $\gamma$ termina in
    $dn^2$ passi.
    Per copiare la prima lettera
in fondo alla parola dopo
il primo blank si compiono $n+1$ passi (si scrive blank nella prima cella e ci si posiziona nella cella $n+2$),
poi se ne compiono $2$ per posizionarsi
sull'ultima lettera della parola originale,
ulteriori $n-1$ per posizionarsi sul primo blank
che si incontra (che corrisponde alla posizione della prima lettera della parola $w$)
e infine $1$ per riposizionare la lettera
cancellata all'inizio e tornare in $q_0$.
Per un totale di
\[
    (n+1)+ 2 + (n-1) + 1 = 2n+3
\]
passi.
Per copiare la seconda lettera,
si compiono
$n-1$ passi per arrivare al primo blank
e poi altri $2$ passi per posizionarsi
nella prima casella vuota dopo la parola già copiata.
A questo punto ci vorranno $3$ passi
per scrivere la lettera nuova e posizionarsi
sull'ultima lettera della parola originale,
ulteriori $n-2$ per posizionarsi sul primo blank
che si incontra (che corrisponde alla posizione della seconda lettera della parola $w$)
e infine $1$ per riposizionare la lettera
cancellata all'inizio e tornare in $q_0$.
Per un totale di
\[
    (n-1)+ 2 + 3+ (n-2) + 1=2n+3
\]
passi.
Facendo questo ragionamento per ogni lettera della parola
si trova che il costo non dipende dalla lettera,
quindi la macchina scrive sul nastro la parola
$w\square w$ dopo $n(2n+3)$ passi e si trova sul
blank tra le due parole.
    A questo punto con $3$ passi
    la macchina scala una lettera indietro, e quindi dopo $3n$
    passi siamo nella situazione $ww$, e alla macchina 
    serviranno solamente altri $2$ passi per arrestarsi.
    Quindi per traslare la seconda parola si compieranno
    \[
    3n+2
    \]
    passi.
    In totale, la macchina compierà 
    \[
    n(2n+3) +3n+2
    \]
    passi e si posizionerà sul primo blank dopo la parola.
    In conclusione abbiamo che $\mu$ decide $X''$ in tempo
    \[
        T'(n):=bn^2 + (n+2) + n(2n+3) +3n+2 + (2n+2) + c^2n^2
    \]
    per un'opportuna costante $c'$ abbiamo chiaramente che
    \[
    T'(n) \leq c'n^2
    \]
    Abbiamo mostrato quindi che
    $X'' \in \operatorname{DTIME}_1^A(c'n^2)$.
\end{sol}